\section{Relations among periods on a theta chain}\label{sec:transport}

In the slope estimate a period contributes both a complex norm and finite-place divisibilities. To compare them, follow the same period as the finite data in its theta kernels vary. We fix the archimedean data for the chain and vary the admissible finite
tuples, including the reference tuple. Unfolding expresses the
resulting pairings through common class sums. The finite
valuation chapter uses this identity to estimate the same
period maps at the places where their lattices differ.
Lemma~\ref{hyp:joint-period-towers} identifies the fixed tower and its
Haar measures. In the descending theta integral \eqref{eq:down-theta},
the subscript \(t\) records the finite tuple on this fixed tower.

\begin{established}[Common period data along a word]\label{hyp:joint-period-towers}
For one word \(u\), fix its reference tuple \(0\) and write \(m=2s+2\)
for the even chain length. Write \(V_j,M_j^\infty\) for the descending
space and its smooth representation at infinity, \(X_j\) for its
ascending definite companion, \(F\) for the maximal totally real
subfield, and \([U]=U(F)\backslash U(\mathbb A_F)\). At all tuples
with this word the measures, additive character and real Weil actions agree,
and the last pair of spaces is
\[
 V_m=T,\qquad X_m=T\perp L_z,\qquad [U(T)]\text{ compact}.
\]
\end{established}
\begin{proof}
The tuple construction in Proposition~\ref{prop:translate-completion}
changes the indicated finite characters and test vectors on a fixed
Hermitian tower. The real spaces, orthogonal complement, additive character
and Haar normalizations are fixed before that change.
For \(G=U(T)\) take the integral compact open \(K_f\subset G(\mathbb A_{F,f})\).
The class-number theorem gives a finite double quotient
\(G(F)\backslash G(\mathbb A_F)/(G(F_\infty)K_f)\)
\cite[Theorem~5.1]{platonov-rapinchuk-class}. Since \(T\) is a subspace
of the definite terminal companion, \(G(F_\infty)\) is compact.
The adelic quotient is a finite union of continuous images of
the compact group \(G(F_\infty)K_f\).
\end{proof}

A column \(\varphi\) at level \(j\) gives scalar automorphic
functions \(r_{t,j,\alpha}(\varphi,y)\) on \(U(X_j)\), one for
each finite chart \(\alpha\) and frame \(y\) at infinity. A row
\(G\) gives a cusp form \(G(v)\) on \(U(V_j)\) for each
\(v\in M_j^\infty\). If \(e=\sum e_g[g]\) has finite support at
the finite adeles, set \(R(e)f(x)=\sum_g e_g f(xg)\). Write
\(B_{t,j}(\varphi,G;y,v)\) for
\[
 \sum_\alpha {\cal B}_{j,\rho}
   \bigl(r_{t,j,\alpha}(\varphi,y),R(e_{t,j,\alpha})G(v)\bigr),
 \qquad j+2\rho+2=m .
\]
Here \({\cal B}_{j,\rho}\) integrates on the rational quotient
of the Bessel group, the semidirect product of the flag
unipotent by \(U(X_j)\), with its specified character on the
unipotent. The top functional is the compact period
\[
 \sum_\beta\int_{[U(T)]}
    R(e_{t,m,\beta})G(v)(h)\,r_{t,m,\beta}(\varphi,y)(h\oplus 1)\,dh .
\]
The finite pairing remaining after the archimedean branching map is
denoted by \(b_t(\varphi,G)\).

Denote the one-step column map by \(K_{t,j}\), the level
endomorphism by \(S_{t,j}\), and the descending row map by
\(D_{t,j}\), with terminal argument on \(T\) for \(j=m-1\).
Write \(I_j(P,v)\in M_j^\infty\) for the common archimedean
intertwiner, where \(P\) is a Fock polynomial of the appropriate
matrix size. With the theta kernel \(\Theta_{t,j,P}\) and a
descending form \(G\) the normalization of \(D_{t,j}\) is
\begin{equation}\label{eq:down-theta}
 \int_{[U(V_{j+1})]} \Theta_{t,j,P}(g,h)G(v)(h)\,dh
     = i_{t,j}\,(D_{t,j}G)(I_j(P,v))(g).
\end{equation}
The last integral uses \([U(T)]\). Each \(i_{t,j}\) is nonzero.
On the designated, ramified column subspaces both \(K\) and \(S\)
preserve the required level and type. Define the double step and
the two scalar products by
\begin{align}
 A_{t,j}&=K_{t,j+1}K_{t,j}S_{t,j}, \qquad
 U_t=A_{t,m-2}\cdots A_{t,0} \quad(j\ {\rm even}
                       \text{ in the product}),\nonumber\\
 h_t&=p_t\prod_{a=0}^s v_{t,2a}\,k_{t,2a}\,k_{t,2a+1},
 &d_t&=\prod_{j=0}^{m-1}i_{t,j}. \label{eq:top-bottom}
\end{align}
Here \(k\) normalizes a single definite lift, \(v\) the two-step
period unfolding, and \(p\) the terminal compact pairing.

\subsection{An identity for definite lifts}

To descend a relation involving several tuples, expand the compact real average in frames that work simultaneously for every term of the relation.
At each real place the ranks \(a,b\) of consecutive definite
spaces satisfy \(0<a\le b\). At a fixed finite lattice class, the column is represented by a polynomial \(F\) of the prescribed
determinant-twisted highest weight in an \(a\)-frame, evaluated
at \(y\), with its prescribed determinant factor. The compact
real group in this basis maps onto \(\prod_\tau U(a)\). The
lifting kernel at a pair of finite classes defines a linear
functional \(\Lambda\) on the Fock polynomials on rectangular
matrices. Both \(F\) and \(\Lambda\) may depend on the tuple. Write
\(\Theta^X\) for the integral operator on the definite quotient.

\begin{lemma}[A common frame expansion]\label{lem:common-definite-expansion}
For the stated determinant-twisted polynomial type and a fixed
real test \((P,y)\), there are finitely many frames \(y_\ell\) and
coefficients \(a_\ell\) independent of the tuple, the finite target
chart and the column such that
\begin{equation}\label{eq:uniform-frame}
 \sum_\alpha \Theta^{X}_{t,j,\alpha\gamma,P}
                   r_{t,j,\alpha}(\varphi,y)
   = k_{t,j}\sum_\ell a_\ell\,
                r_{t,j+1,\gamma}(K_{t,j}\varphi,y_\ell)
\end{equation}
for every \(t,\gamma,\varphi\) of this type. For a two-step lift the
expansion can be iterated with coefficients common to the tuples
and with the remaining scalar \(k_{t,j}k_{t,j+1}\).
\end{lemma}

\begin{proof}
\proofstep{Isolate the compact frame average.}
The average to be computed, apart from determinant factors, is
\[
 \int_{\prod_\tau U(a)}
  \Lambda(u^{-1}\cdot P)\,F(uy)\,du ,
 \qquad \int du=1.
\]
Split the definite quotient into its finite compact-level classes.
Each integral is its compact real average times the finite level
mass divided by the order of its rational stabilizer (with the
real Haar mass restored if a probability is used).
The two compact-level characters in the integrand cancel.
These are precisely the weights of the finite class lift.

\proofstep{Identify the equivariant multiplicity space.}
Before applying \(\Lambda\), take the compact-group average
\[
 Q(F)=\int F(uy)(u^{-1}\cdot P)\,du .
\]
The polynomial orbit has bounded multidegree, so this integral
is defined coefficient by coefficient. Left translation in Haar
measure makes \(Q\) equivariant for the small compact group.
To determine its image, first
work at one place. The polynomial algebra on \(b\)-by-\(a\)
matrices has the Cauchy decomposition into
\(S_\lambda(\C^a)\otimes S_\lambda(\C^b)\), with the duals
determined by the action on variables. The multiplicity space
of the specified irreducible \(U(a)\) polynomial type is thus
irreducible for \(\mathrm{GL}_b(\C)\). Substitution of the first
\(a\) rows defines a nonzero intertwiner, since arbitrary square
matrices occur there. Its translates by
\(z\in\mathrm{GL}_b(\C)\) span the multiplicity space.

For a product of places the equivariant homomorphism space on
each homogeneous component is the tensor product of these
multiplicity spaces, by applying the decomposition at each
place in turn. Consequently
\[
 Q(F)=\sum_{\ell=1}^h s_\ell\,\operatorname{sub}_{z_\ell}F .
\]
Each \(z_\ell\) is an invertible frame at every place. The list
depends on \(P,y\) and the common representations. It works
simultaneously for all \(\Lambda,F\). Restoring the determinant
factors multiplies \(s_\ell\) by the source factor at \(y\) and
the target factor at the scaled frame \(y_\ell\). This scaling is
common to all finite classes. The defining finite class matrix
therefore turns the frame expansion into \eqref{eq:uniform-frame}.
Independence from the finite chart lets one expand the whole linear relation in the same frames.

\proofstep{Expand the two oscillator variables.}
A two-step unfolding gives a smooth function \(C(x,x')\) on
two successive oscillator spaces with the finite self-dual lattice
vacuum in each variable. It suffices to make a finite polynomial
expansion after projection to the central characters occurring in
the definite lifts. For a compact centre \(Z\), with Haar probability measure and
character \(\omega\), this projector acts on a function \(F\) by
\[
        P_\omega F=\int_Z\omega(z)^{-1}\rho(z)F\,dz ,
\]
where \(\rho\) is the given central action.
Moving the central variable around the integral shows that
inserting this projection preserves the lift. The centre acting
on the intermediate group acts on both kernels, so its matching
character on the second factor is obtained by this same change
of integration variable. These projections work for the whole
family because the real Weil actions and the central characters
of the readings agree.

In the definite Fock space a central circle
acts on each polynomial degree by that degree shifted by the
vacuum character (with sign from the polarization). More explicitly, if the vacuum has central exponent \(b\) and
the polarization has sign \(\varepsilon\in\{1,-1\}\), a homogeneous
slice of degree \(d\) and the projection onto exponent \(a\) contribute
\[
 \int_{|z|=1} z^{b+\varepsilon d-a}\,dz
   =\mathbf1_{\{a=b+\varepsilon d\}},
\]
with probability Haar measure. At each place only the prescribed
nonnegative degree can survive. The
slices thus lie in fixed finite-dimensional polynomial spaces
times their vacua. Denote the projected function again by \(C\).

For the first space choose evaluations at \(x_1,\ldots,x_h\)
forming a dual basis: their joint annihilator is zero. Let
\(L_a\) be the interpolation basis, so \(L_a(x_b)=\delta_{ab}\).
For every \(x,x'\),
\[
 C(x,x')=\sum_{a=1}^h L_a(x)\,C(x_a,x').
\]
Each \(C(x_a,\cdot)\) belongs to the second slice space.
Apply \eqref{eq:uniform-frame} in the first and then the second
variable of each summand. At a fixed final frame the coefficient
is the sum of products of the two transition coefficients over
intermediate frames. Only \(k_{t,j}k_{t,j+1}\) varies with the
tuple outside the final column reading.
\end{proof}

\subsection{Unfolding and the finite measures}

The common frame expansion leaves a scalar at each tuple. Its finite factors come from the orbit surviving unfolding and the measures of its level cosets. At the terminal step
it uses the compact branching projection.

\begin{lemma}[The orbit surviving the flag period]\label{lem:joint-flag-orbit}
In the mixed Schr\"odinger model with fixed frame \(x_0\), the surviving rational frames \(x\) form a single orbit characterized by
\[
 \operatorname{rank}x=\operatorname{rank}x_0,\qquad x^*Hx=x_0^*Hx_0.
\]
 The remaining integral is over
the stabilizer quotient and converges absolutely.
\end{lemma}
\begin{proof}
Expand the theta kernel in rational matrices \(x\) with \(\rho\)
columns. On the rational quotient of the flag centre the normalized
integral of the character with parameter \(x^*Hx-x_0^*Hx_0\) is
\[
 \begin{cases}1,&x^*Hx=x_0^*Hx_0,\\0,&\text{otherwise}.\end{cases}
\]
For a dependent surviving matrix choose \(a\) with \(xa=0\), its
last nonzero entry \(a_j=1\). The flag relation supplies a
transvection fixing this term on the first \(j\) columns, with
nontrivial flag character \(\xi(u)\). Translation invariance of the mean gives
\[
 (1-\xi(u))I_x=0,\qquad \xi(u)\ne1 .
\]
Thus the dependent term has zero integral. This exclusion also applies
when the prescribed Gram matrix is zero.

For full-rank \(x\), the map sending the columns of \(x_0\) to
those of \(x\) is an isometry on their span. Witt extension makes
it the restriction of a rational unitary transformation. In the
product with the rational flag group the stabilizer consists of
pairs satisfying \(g x_0=x_0u\), with \(u\) upper unitriangular.
The sum over its rational cosets unfolds to the stabilizer quotient.
The absolute theta sum for a fixed Schwartz function grows at most
as a power of the height. The flag unipotent quotient is compact,
and rapid decay of the cusp form makes this majorant integrable.
\end{proof}

\begin{lemma}[Reduction to the next period]
After the two stabilizer integrations, the integrand is the translated period at \(j+2\); its two remaining variables range over
\[
 R_1(\mathbb A_F)\backslash U(V_{j+1})(\mathbb A_F)\quad\text{and}\quad R_2(\mathbb A_F)\backslash U(V_{j+2})(\mathbb A_F).
\]

\end{lemma}
\begin{proof}
For the Bessel-to-Fourier--Jacobi step the intermediate stabilizer
is the Jacobi group \(R_1\). The mixed-model function evaluated
on its base frame is the lower theta function, with the prescribed
Bessel line in the last coordinate. In the Fourier--Jacobi-to-Bessel
step it is paired along the line coordinate before restriction
to the next Bessel stabilizer \(R_2\). In each step write the
stabilizer as \(N\rtimes U_0\). The measure is the image of
\(dn\,du_0\) under \((n,u_0)\mapsto n u_0\), with the rational
unipotent quotient of mass one.

Since the rational points split as \(N(F)\rtimes U_0(F)\), the
quotient formula for an integrable pairing \(F_0\) is
\[
 \int_{[N\rtimes U_0]} F_0(h)\,dh
    =\int_{[U_0]}\int_{[N]}F_0(nu)\,dn\,du .
\] The former integral evaluates the flag character; the
latter is the remaining Bessel or Jacobi pairing. Apply the quotient
formula at \(R_1\) and \(R_2\) successively. Absolute convergence
from Lemma~\ref{lem:joint-flag-orbit}, uniformly on compact representatives for the definite
factor, permits both integrations.
\end{proof}

\begin{lemma}[Mass away from the twist places]
Assume the fixed-level support conditions of
Section~\ref{subsec:fixed-level-support}. Off the twist places \(Q_t\) let \(K_i^*\) be the product of the
prescribed fixed levels and the full lattice stabilizers elsewhere.
The two integrations contribute a common scalar \(\kappa_j=v_1^*v_2^*\),
where
\[
 v_i^*=\frac{\mu_i(K_i^*)}{\rho_i(R_i\cap K_i^*)}.
\]
Here \(\mu_i,\rho_i\) are the ambient and stabilizer measures.
\end{lemma}
\begin{proof}
Outside the twist and fixed level places the vectors are the unramified lattice vacua.
The test for the unfolded function to be nonzero places the
group variable in \(R_iK_i\). At the fixed level places the
corresponding test after the prescribed Heisenberg translates
places it in \(R_iJ_i\), with \(J_i\) the specified compact
level. The transformed integrand is left invariant by \(R_i\)
and right invariant by these levels. In the first variable the
two intermediate level characters cancel each other. In the
second the row has the indicated level character. The second
support test is used after the first variable has been reduced.

For the indicator of \(K_i^*\) the fibre integral equals
\(\rho_i(R_i\cap K_i^*)\) on the image of \(K_i^*\) in the quotient,
and is zero elsewhere. Hence the quotient formula reads
\[
            \mu_i(K_i^*)=v_i^*\,\rho_i(R_i\cap K_i^*),
\]
which gives the asserted mass. At \(Q_t\) normalize the
maximal compact and its intersection with \(R_i\) to mass one.
Factor the finite adele and stabilizer measures at \(Q_t\) and
its complement. The factors outside \(Q_t\) contribute the two
masses just computed, leaving the twist integrals.
\end{proof}

\paragraph{Test functions at a split twist place.}

The two finite test functions determine the cell weights in the unfolding. Let \(k_0\) be the local field and
\(K_i^0\) the full integral general linear group in the two
intermediate variables. In the local matrix model the additive
character has conductor \(\mathcal O_{k_0}\) and the additive
measure of \(\mathcal O_{k_0}\) is one. For a lattice \(L\) in the variables being transformed let
\(L^\perp\) be its annihilator for the given pairing and character.
Translation by any \(a\in L\) multiplies the Fourier integral by
\(\psi(\langle a,y\rangle)\). Hence
\[
 \widehat{\mathbf1_L}(y)=\int_L\psi(\langle x,y\rangle)\,dx
       =\operatorname{vol}(L)\mathbf1_{L^\perp}(y).
\]
The standard integral coordinate lattice has volume one. The two local vacua can be read as these
indicators, followed by the finite Heisenberg combinations
defining the test functions \(P_1,P_2\). Transform each to
position one (the flag coordinate carrying the additive
character).

In the source this operation on a level vector
at a representative \(\alpha\) is expressed as a finite
combination \(e_\alpha\) of group translations. For every equivariant functional \(l\) on \(\tau\) and
level-character vector \(w\), the finite expansion gives
\[
 l\big(\tau(e_\alpha)w\big)=\sum_a e_\alpha(a)\,l\big(\tau(a)w\big).
\]
Thus the weights used on each cell can be evaluated before summation. The same formula applies
to the target vector.

Choose representatives \(\gamma,\beta\) for the successive
nonzero cells of the two transformed functions. The test-data
conditions here are that the representatives cover the support
and give distinct \(R_i\)-level double cosets in each variable.
Let \(\xi_1(\alpha,\gamma)\) and \(\xi_2(\gamma,\beta)\) be the Schwartz
functions on these cells given by partial Fourier restriction and
unipotent averaging. At compact open level \(J_i\subset K_i^0\) the
stabilizer in the position-one group \(R_i\) is
\[
                      C_i(x)=R_i\cap xJ_ix^{-1}.
\] Write
\(\pi_i\) for projection onto the smaller general linear factor
and \(\nu_i\) for Haar measure on it.

\begin{lemma}[The contribution of the twist cells]\label{lem:joint-twist-cells}
At a split twist place, use the two position-one test functions.
Assume their nonzero cells exhaust the support and the representatives
\(\gamma,\beta\) give distinct stabilizer--level double cosets in
each variable.
If \(\mathcal T\) is the remaining period, trilinear in the two
Schwartz functions and the row, the local integral is
\begin{multline}\label{eq:twist-cell}
 \frac{1}{[K_1^0:J_1][K_2^0:J_2]}
 \sum_{\gamma,\beta}
 \frac{1}{\nu_1(\pi_1 C_1(\gamma))\nu_2(\pi_2 C_2(\beta))}
 \\ {}\times
 \sum_b e_\beta(b)\,
 \mathcal T\bigl(\xi_1(\alpha,\gamma),\xi_2(\gamma,\beta),
                           \tau(b)w\bigr).
\end{multline}
\end{lemma}
\begin{proof}
Disintegrate over \(R_i\backslash R_i xJ_i\).
The numerator of the cell mass is \(\mu_i(J_i)=1/[K_i^0:J_i]\).
In the stabilizer integral first integrate
in the radical coordinates with the self-dual additive measure.
More explicitly, write \(N_i\) for the radical and \(m_{N_i}\)
for its measure in position one. With \({\rm flat}_i\) the
position-one partial Fourier restriction of the translated
test function, the entry used here is
\[
 \xi_i(x',x)= {\rm flat}_i(x',x)\,
   \frac{m_{N_i}(N_i\cap K_i^0)}
        {m_{N_i}(N_i\cap C_i(x))}.
\]

The denominator records the actual radical volume at the cell
and the numerator uses the reference lattice. The remaining
stabilizer measure is \(\nu_i(\pi_i C_i(x))\). Right
translation by a level element multiplies \(P_i\) and the
vector or adjacent kernel by cancelling characters, so the
cell integral has the displayed value for any representative.
Applying this identity in the inner and then the outer variable
gives the product of the two cell masses. Substitution of the
finite translate \(e_\beta\) and trilinearity give \eqref{eq:twist-cell}.
\end{proof}

\begin{proposition}[Scalar of the two-step period integral]\label{prop:joint-scalar-descent}
For the data of Lemma~\ref{hyp:joint-period-towers}, with the local cell
decomposition of Lemma~\ref{lem:joint-twist-cells}, the scalar of
the unfolded integral at a nonterminal step is
\[
 v_{t,j}=\kappa_j\,s_{t,j}s_{t,j+1}c_{Q,t,j}c_{G,j}c_{T,t,j}.
\]
The terminal step has \(\kappa_j=1\) and its compact pairing
contributes \(p_tb_t\).
\end{proposition}
\begin{proof}
At several twist places use product representatives in \eqref{eq:twist-cell}.
Distinct places act on distinct local
matrix coordinates. For distinct places \(v,w\) the local operators \(A_v,A_w\)
commute on the tensor product. If \(b_v(t_v)\) denotes a local
weighted cell entry in \eqref{eq:twist-cell}, the corresponding
global sum consequently factors as
\[
 \sum_{(t_v)_v}\prod_v b_v(t_v)=\prod_v\sum_{t_v} b_v(t_v).
\]
The same commutation holds with modifications at disjoint fixed
level places. Thus the full twist-place sum gives precisely the chart
sum of the two definite lifts, with their entry weights in
\eqref{eq:twist-cell}, and an overall scalar \(c_{Q,t,j}\).
The real datum after contraction can be read at any reference
tuple. Restriction to the frame and Fourier integration use only
the common splitting and measures after the finite translates have
been removed.

By bilinearity, replacing the two test functions by their
multiples changes their contraction \(B\) by
\[
 B(s_{t,j}\phi_1,s_{t,j+1}\phi_2)
                 =s_{t,j}s_{t,j+1}B(\phi_1,\phi_2).
\]
\(c_{G,j}\) is the common fixed-level factor and \(c_{T,t,j}\)
the twist of the other prescribed local factors. At the last
step the compact terminal period is already an integral on
\([U(T)]\), and the same formula uses \(\kappa_j=1\).

For that compact period write \(U(T)\) in finite level classes.
Its real integral is the equivariant projection for the
restriction from \(U(X_m)\) to \(U(T)\) along \(h\mapsto
h\oplus1\). Let \(V\) and \(W\) be the selected terminal types on \(U(X_m)\)
and \(U(T)\), respectively. Their highest weights interlace, so
unitary branching gives
\[
 \dim\operatorname{Hom}_{U(T)}(V|_{U(T)},W)=1 .
\]
For fixed \(y,v\) the real integral is therefore a scalar multiple
of the prescribed branching map. The scalar is common for all \(t\),
and summing the finite classes leaves \(p_tb_t\).
\end{proof}

\subsection{Support at the fixed level places}
\label{subsec:fixed-level-support}

At the fixed level places the scalar in the unfolded integral is determined by the support of the two translated test functions. At a fixed level place the first flat map
has support in the product of the Jacobi subgroup and the
specified level. After contraction, the second flat map
has support in the product of the Bessel subgroup and its
specified level. These statements concern the translated
test functions on the chosen level lattices. We prove them
for the split model of the fixed level. In the formulas below
the conductor of the additive character \(\psi\) is \(O\),
the integer ring, and its additive Haar measure \(\nu\)
gives \(O\) mass one.

\begin{lemma}[Support of the split position functional]\label{lem:split-position-support}
At a fixed split level place with the symmetric chain and
its standard iterated type function \(\Phi\), let \(c\) be the number of numbered target coordinates. For \(1\le p\le c\), let \(P_p\) integrate the test function against the prescribed additive character over the free entries of the numbered slice at \(p\). With \(R_p\) the position subgroup and \(J\) the prescribed level, the support satisfies
\[
 \{h:P_p(h\Phi)\ne0\}\subseteq R_pJ.
\]

\end{lemma}

\begin{proof}
\proofstep{Record the integral flag coordinates and their residue support.}
Write \(k\) for the local field, \(\varpi\) for a uniformizer,
and \(\Phi\) for the standard iterated type function at a stage
of a tuple. There are \(c\) numbered target flag coordinates
and \(c+1\) source flag coordinates, with ordinary coordinates
in the remaining block. With the ordinary block first, the target and source frame
matrices have the form
\[
 D=I_{\rm ord}\oplus\operatorname{diag}(\varpi,\ldots,\varpi^c),
 \qquad D'=I_{\rm ord}'\oplus
             \operatorname{diag}(\varpi,\ldots,\varpi^{c+1}).
\] The field coordinates of
the test matrix \(X\) are read as \(D^{-1}XD'\). In these
coordinates the level \(J\) is integral. The entries on
and below the prescribed parts of the flag diagonal agree
with the identity modulo \(\varpi\). The iterated test keeps
the retained columns of a matrix in that residue pattern;
its other coordinates are Schwartz coordinates. Let \(M\) be the corresponding rectangular matrix space. In the
original coordinates the residue condition reads
\[
 \operatorname{supp}\Phi\subseteq D M(O)D'^{-1},\qquad
 \Phi(X+DZD'^{-1})=\Phi(X)\quad(Z\in\varpi M(O)).
\]

At a numbered position \(1\le p\le c\), the operator \(P_p\)
integrates \(\Phi(X)\psi(-\sum_i X_{i,i+1})\) over the free
entries of \(X\). In the character sum the two indices run
in the numbered flag blocks. The ordinary block is the
argument of \(P_p\Phi\), and the numbered diagonal entries
of the affine slice are one.

Write \(R_p\) for the group
generated by the ordinary general linear block and the
target unipotent for this slice. The latter has scalar
covariance under \(P_p\), and the ordinary block acts by
invertible changes of the argument and a nonzero scalar.
Let \(m(h)\) be the product of the determinant character
and modulus for a target translation. Its Weil action satisfies
\[
 (h\Phi)(X)=m(h)\Phi(h^{-1}X),\qquad m(h)\ne0.
\]
In particular the scalar does not alter the translated support. These two covariance
statements follow by changing the free variables. They
apply to the same numbered slice before and after translation.
We claim that
\begin{equation}\label{eq:position-level-support}
                   P_p(h\Phi)\ne0\quad\Longrightarrow\quad h\in R_pJ .
\end{equation}

For the first reduction use Iwasawa on \(hD\), with weights
\(2p-1\) on the ordinary block and \(2i\) on numbered
line \(i\). This writes \(h=n M w D^{-1}\), where \(w\) is
integral with integral inverse, \(n\) is in the target
unipotent of \(R_p\), and \(A=M^{-1}\) has block form
\[
          A=\begin{pmatrix}P&Q\\0&S\end{pmatrix},
       \qquad S=\operatorname{diag}(s_i),\qquad Q_i=0\ \ (i\ne p).
\]
Here \(Q_i\) is column \(i\) of \(Q\). This form follows by
taking the Levi for the order of the numbered lines with
the ordinary block tied to line \(p\), and using the finer
Iwasawa order inside that grade.

Removing \(n\) preserves
nonvanishing. The integral change \(w\in\mathrm{GL}(O)\) preserves both
\(M(O)\) and \(\varpi M(O)\). Thus the transformed coordinates
give the necessary support and the translation lattice
\[
 AXD'\in M(O),\qquad
 Z\in A^{-1}\varpi M(O)(D')^{-1}.
\] At a nonzero reading the
numbered diagonal entries give \(s_i\varpi^i\in O\); column
\(p\) has zero ordinary part and gives \(Q_p\varpi^p\in O\).

For the inverse diagonal condition, let \(y\) range over
\(s_i^{-1}\varpi^{-i}O\). Translate the integration variables
by \(Z=A^{-1}E\), with \(E\) supported at numbered position
\((i,i+1)\) with entry \(ys_i\). Then \(AZD'\) is divisible
by \(\varpi\). The translated variables are free, including
the possible ordinary part of column \(p+1\) when \(i=p\).
The character is multiplied by \(\psi(-y)\).

If the reading is \(I\ne0\), translation invariance gives
\[
 I=\psi(-y)I\quad(y\in s_i^{-1}\varpi^{-i}O)
 \quad\Longrightarrow\quad s_i^{-1}\varpi^{-i}O\subseteq O .
\]
Here \(O\) is the conductor of \(\psi\). Together with the numbered
diagonal condition this proves \(s_i\varpi^i\in O^\times\). With
\(\gamma=P^{-1}\) in the ordinary block,
\(A\operatorname{diag}(\gamma,1)D\) is integral with
integral inverse: its diagonal blocks are the identity and
a diagonal unit matrix, and its upper block is
\(Q\operatorname{diag}(\varpi^i)\).
This gives the factorization required for the first reduction:
\[
 h=r u,\qquad r\in R_p,\qquad u\in D\mathrm{GL}(O)D^{-1}.
\]

For the second reduction it suffices to treat this \(u\), using
the covariance of \(P_p\) under \(R_p\). Nonvanishing gives an
affine-slice value of the translated test. In frame
coordinates this gives a \(\kappa\) in the prescribed residual
level such that \(g=D^{-1}uD\,\kappa\) agrees with the scaled
affine matrix on the retained columns modulo \(\varpi\).
Write \(G=D^{-1}uD\) and let \(Z\) consist of the retained
columns of the scaled affine matrix. The nonzero reading of the
previous type says
\[
             G^{-1}Z\equiv\binom{g_0}{0}\pmod{\varpi},
\]
where \(g_0\) belongs to the previous level. In the reordered
frame \(\kappa=\operatorname{diag}(g_0,1)\) is in the next level.
Thus the retained columns of \(G\kappa\) agree with \(Z\) modulo
\(\varpi\).
The affine zeros imply the following two
conditions. The lower entries in the retained initial
ordinary columns, in rows \(i\ge p\), vanish modulo
\(\varpi\).

In every numbered column the entries at or below
that column's grade agree with the identity modulo
\(\varpi\). Replace those entries by their exact zero and
one values, fixing all other entries. Writing \(\widetilde g\) for the changed matrix, one has
\(g^{-1}\widetilde g\equiv1\pmod\varpi\), so this change is made
inside the prescribed right level. Its lower numbered square
block \(S_0\) is upper triangular unipotent. If \(E_0\) is
the lower rectangular block, then
\[
     (E_0\ S_0)
         \begin{pmatrix}1&0\\-S_0^{-1}E_0&1\end{pmatrix}
          =(0\ S_0).
\]
The multiplying block is zero on the retained initial ordinary
columns because \(E_0\) is zero there. On the other columns this
shear belongs to \(J\) by the iterated level conditions.
The result belongs to the frame reading of \(R_p\) up to
principal congruence. Each right factor used belongs to
the prescribed level. This proves \eqref{eq:position-level-support}
with that level rather than the full lattice stabilizer.
\label{cor:period-level-support-1}

\paragraph{Exchange of adjacent positions.}
The position calculation uses the symmetric chain of even flag length, so the midpoint and full-length parameters satisfy
\[
 s=2\rho,\qquad 1\le i\le\rho\quad\text{in the first half}.
\] For a move from position \(i\) to \(i+1\) with
\(i\le\rho\) and \(i+1\le c\), the two affine integrals have the
same free variables except for an ordinary row and an ordinary
column. The integral for \(i+1\) frees the entries in numbered
target row \(i\) and the ordinary source columns. The integral
for \(i\) instead frees the entries in numbered source
column \(i+1\) and the ordinary target rows. The compensating blocks lie in the kernel of the additive character in the two affine integrals: if \(b\) denotes the change in its argument in either block,
\[
 \psi(a+b)=\psi(a)\psi(b)=\psi(a).
\]

The first block contributes
the indicator of \(\varpi^iO\) in each coordinate: the frame
divides these entries by \(\varpi^i\), and the permitted
integral row addition preserves the test. This addition has
unipotent determinant one and leaves the residual type
character trivial on the added row. The second contributes
the indicator of \(\varpi^{-i}O\): the frame multiplies the
column by \(\varpi^{i+1}\), and these entries are prescribed
zero modulo \(\varpi\). In each free entry the two tests reduce respectively to
\[
 \varpi^{-i}x\in O\ \Longleftrightarrow\ x\in\varpi^i O,
 \qquad \varpi^{i+1}y\in\varpi O\ \Longleftrightarrow\
                       y\in\varpi^{-i}O .
\]
After removing these indicators both restrictions give the same
test with the free block zero.

For the row this uses the iterated row pattern, for which
position \(i\le\rho\) permits the shear after the initial
row has been set to zero. For the column it uses the
periodicity modulo \(\varpi\) of the frame reading.
Thus the ratio of the two operators is the positive scalar
\[
     \frac{\nu(\varpi^iO)^{d_{\mathrm{src}}}}
          {\nu(\varpi^{-i}O)^{d_{\mathrm{tgt}}}},
\]
where the exponents are the respective ordinary block
dimensions. Iteration compares the operator at any
position needed here to the operator at position one.
Empty ordinary blocks contribute factors one. The ratio
depends on the place, dimensions and position and is common
to the tuples, which use the same lattice exponents and
the same measure.
\end{proof}

\paragraph{The two flat maps.}
Use Lemma~\ref{lem:split-position-support} at each position of
the symmetric chain.
In the Bessel to Fourier--Jacobi step the flat column map
at position \(p\) with extra column \(y\) reads
\(P_p(E_y^{-1}h\Phi)\), with \(E_y\) the upper shear that
adds the extra column into the ordinary block and has
determinant one. This follows by translating the affine
matrix by the added column before taking the Fourier
integral. Equation~\eqref{eq:position-level-support} yields
\[
  P_p(E_y^{-1}h\Phi)\ne0
     \ \Longrightarrow\ E_y^{-1}h\in R_pJ
     \ \Longrightarrow\ h\in E_yR_pJ .
\]
In the split frame the product \(E_y r\) lifts to the
Jacobi group. The right factor lifts to precisely the
target level at that stage because the other split block
is inverse adjoint. This gives the first support statement.

For the value calculation at the identity, the middle
column and row give pure factors. With \(c=2a\) in a double
step, the first flat column in its middle variable is
a position-one entry times a positive constant times
\(\mathbf1_{\varpi^{-a-1}O^d}\). The following flat row is
its own position-one entry times a positive constant times
\(\mathbf1_{\varpi^{a+1}O^d}\), with \(d\) the middle
ordinary dimension. For the added entries \(y\) in numbered column \(a+2\), the
source frame and its zero residue condition give
\[
  \varpi^{a+2}y\equiv0\pmod{\varpi O^d}
         \quad\Longleftrightarrow\quad y\in\varpi^{-a-1}O^d .
\]

In this
range the addition is a translation by the period lattice
of the test function.
For the row the ordinary-to-numbered conversion is the
source shear with the direct range \(\varpi^{a+1}O^d\).
It has determinant one and belongs to the source level
with trivial character at these row coordinates. Outside
the range, the affine row test is zero by the dual lattice
condition. The remaining operator in each case is put at
position one by the exchange calculation. The position-one
normalization restores the positive Haar volume of its
radical stabilizer.

To compare the two position-one entries, let \(C_0,R_0\) denote those entries and \(c_C,c_R>0\) their pure-map factors. The common lattice is
\[
 L=\varpi^{-a-1}O^d .
\] The column and row functions are
\(c_C C_0\mathbf1_L\) and \(c_R R_0\mathbf1_{L^\perp}\).
For the additive character \(\psi\) of the middle model,
\[
 \widehat{\mathbf1_L}(z)=\int_L\psi(\langle x,z\rangle)\,dx
             =\nu^d(L)\mathbf1_{L^\perp}(z),
 \qquad L^\perp=\varpi^{a+1}O^d.
\]
Indeed changing \(x\) by \(x_0\in L\) multiplies the integral
by \(\psi(\langle x_0,z\rangle)\); it is therefore zero unless
this equals one for every \(x_0\). At the identity the
contraction is \(c_Cc_R C_0R_0\nu^d(L)\nu^d(L^\perp)\).

For the translated second test let \(R_h(z)\) be its row
function. The contraction is
\[
 c_C C_0\nu^d(L)\int_{L^\perp}R_h(z)\,dz .
\]
On \(L^\perp\) the source shear belongs to the level and
commutes with the left target translation; thus \(R_h(z)\)
is the position reading of \(h\Phi\) multiplied by its
nonzero shear character and pure-map factor. A nonzero value
forces that reading to be nonzero, whence \(h\in R_pJ\) by
\eqref{eq:position-level-support}. Here \(R_p\) lifts to the
target Bessel subgroup. Only the transformed column confines
this integral. The support of the translated second test can
be larger.

Both deductions use the split Weil
model and the compatible Jacobi and Bessel frames of
the same step. In the second frame let \(e_\xi\) denote the image of the
negative line generator under the fixed line identification. The
normalization in this split model is
\[
 \langle e_\xi,e_\xi\rangle=-\xi^2 .
\]
\label{cor:period-level-support-2}

\paragraph{Use at all the fixed places.}
Let \(S\) denote the fixed places and write the translation
as \(h=\prod_{v\in S} h_v\). The operators on distinct local
Weil factors commute. For each \(v\), let \(T_v\) be the local contraction or position map on the corresponding Weil factor. Applied to a finite expansion \(\sum_a c_a\bigotimes_v\phi_{a,v}\), the tensor product of these maps gives
\[
  \sum_a c_a\bigotimes_v T_v(\phi_{a,v}).
\]

If one \(T_v\) vanishes on the full local test span, each summand
vanishes, independently of the other places.
Thus a nonzero global value implies membership in the local cell at every \(v\), giving
\(h\in\prod_{v\in S} R_vJ_v\) in the corresponding local
notation. Multiplying the local Jacobi or Bessel factors
and the local levels gives the fixed-level cell, using the
specified product level. At the identity the common
multiplier is the product of the positive local constants.
The terminal two steps use the rank-one line identification
and compact terminal pairing for this common-value assertion;
the double-step support computation concerns the
nonterminal steps. These are the cell integration inputs.


\label{cor:period-level-support-3}


\subsection{Descent of a linear relation}

The common frames and the finite measure factors now enter the same induction. With the indicated coefficient change, a relation at the higher level descends through the two lifts.

\begin{construction}
\[
 \begin{array}{rcccl}
  &\text{level }j&&\text{level }j+2&\\[2pt]
 \text{column}&\varphi^{(j)}&\xrightarrow{\ A_{t,j}\ }&
                           \varphi^{(j+2)}&\\[2pt]
 \text{row}&G^{(j)}&\xleftarrow{\ D_{t,j}D_{t,j+1}\ }&
                           G^{(j+2)}&\\[3pt]
 \text{coefficient}&a_{i,j}v_{t,j}k_{t,j}k_{t,j+1}
          &=&a_{i,j+2}i_{t,j}i_{t,j+1}.&
 \end{array}
\]
\caption{One even step in Lemma~\ref{lem:joint-period}, with \(t=t_i\). The column moves up and the row moves down. The bottom equality permits the weighted relation at level \(j+2\) to descend.}\label{fig:even-period-step}
\end{construction}

\begin{auxlemma}[The real tests and the period on each marked arrow]
\label{lem:word-real-tests-period}
At stage \(j\) of the word tower of Lemma~\ref{hyp:joint-period-towers},
let \(K_j\) be the product of real maximal compact subgroups on
\(V_j\), and \(H_j=(M_j^\infty)_{K_j\text{-finite}}\).
Let \(\mathcal F_j\) denote the polynomial Fock space on the rectangular
matrices for \((V_j,V_{j+1})\), with the word's real Weil action.
On each used tuple and at every arrow, including the definite last one,
\[
 H_j=\operatorname{span}_{\mathbb C}
         \{I_j(P,v):P\in\mathcal F_j,\ v\in H_{j+1}\}
             \qquad(0\le j<m).
\]
For fixed column and frame the period at every even stage is a
continuous functional on \(M_j^\infty\). The bound is valid separately
in each real factor of a pure tensor.
\end{auxlemma}
\begin{proof}
Realize the descending spaces \(V_j\) as the marked \(Y_j\) of
Proposition~\ref{prop:isometric-refined-theta-chain}, with \(Y_m=T\).
At each stage write the selected isotypic row, at fixed finite level, as
\[
       \mathcal A_{t,j}^{\rm sel}
          =L_{t,j}\otimes M_j^\infty .
\]
Here \(L_{t,j}\) is the finite-dimensional multiplicity space. In
Proposition~\ref{prop:isometric-refined-theta-chain} the symbol \(M_j\)
denotes this multiplicity space. The real factor here is the exterior
tensor product of the specified irreducible discrete series, with
the finite-dimensional compact representation at \(T\).

The common evaluation tensors in that proposition give a nonzero
coefficient of the real intertwiner on each edge, with the same
tensor for both occurrences of each interior stage. In the row integral \eqref{eq:down-theta}, for \(G\in L_{t,j+1}\)
and \(v\in M_{j+1}^\infty\), the two factors are
\[
 G\otimes v\ \longmapsto\ i_{t,j}\, D_{t,j}G\otimes I_j(P,v).
\]
A change of Haar normalization multiplies the former
identity by the scalar \(i_{t,j}\ne0\). The finite chart index and
the finite twist do not enter \(I_j\).

Write \(W_j\) for the span of real images in the statement.
Polynomials have finite compact orbit span, so their images with
compact-finite upper vectors lie in \(H_j\).
The nonzero evaluation can be taken on compact-finite inputs. The pairs are all unitary dual pairs over
\(\mathbb C/\mathbb R\), including the pair with compact upper member.
For these pairs the smooth occurrence of the two factors is also an
algebraic occurrence in the Fock space. In terms of the upper and
lower compact subgroups, this gives
\[
 I_j\bigl(\mathcal F_j,(M_{j+1}^\infty)_{K_{j+1}\text{-finite}}\bigr)
                  \ne 0 .
\]

This is the unitary case of Bao and Sun's Proposition~5.1.
\cite{bao-sun-unitary}
To see the density inside the selected row, fix its multiplicity
vector \(G\) and convolve the upper smooth vector \(v\) on
\(K_{j+1}\). Smooth approximate identities converge in every
smooth seminorm. By the smooth Peter--Weyl expansion on the
compact group, the convolution kernels can themselves be approximated
in these seminorms by finite sums of matrix coefficients. Thus there
are compact-finite \(v_n\) with
\[
       G\otimes v_n\ \longrightarrow\ G\otimes v
               \quad\text{in the selected smooth row}.
\]

Every translate in this construction belongs to the same selected
representation. The real intertwining pairing is continuous on the
smooth oscillator and representation spaces.

Equivariance gives, for a lower compact element \(k\) and a
lower Lie algebra element \(X\),
\[
 k I_j(P,v)=I_j(\omega_j(k)P,v),\qquad
 dM_j(X)I_j(P,v)=I_j(d\omega_j(X)P,v).
\]
The infinitesimal Fock action is a sum of products of polynomial
multiplications and differentiations of order at most two, including
the scalar splitting character. Thus \(W_j\) is a nonzero
Harish--Chandra submodule. Each real factor has irreducible
Harish--Chandra module, and their exterior tensor product is
irreducible for the product group. Consequently
\[
  0\ne W_j\subseteq H_j,\quad
  (\mathfrak g_j,K_j)W_j\subseteq W_j
        \quad\Longrightarrow\quad W_j=H_j .
\]
The same algebraic argument on the core arrows occurs in
Lemma~\ref{lem:fitted-core-span}.

For an even \(j<m\) the integrated group is
\(R_j=N_j\rtimes U(X_j)\). The quotient of \(N_j(\mathbb A_F)\)
by \(N_j(F)\) is compact by its filtration with additive quotients.
The group at infinity of \(U(X_j)\) is compact and the finite
class quotient is finite, as in Lemma~\ref{hyp:joint-period-towers}.
Taking compact representatives successively in the two groups
gives compact representatives for \([R_j]\). At the terminal stage
use \([U(T)]\) directly. Each finite translate in the chart sum
therefore restricts the row to a compact set \(C\) in its automorphic
quotient. With the finitely many translates included, integration gives
\[
 |B_{t,j}(\varphi,G;y,v)|\le C_{t,j,\varphi,y}\,
                \sup_{h\in C}|G(v)(h)|.
\]

On a finite cover of that set by real coordinate charts, local
Sobolev embedding for an integer \(d>\dim_{\mathbb R} U(V_j)(F_\infty)/2\)
bounds the supremum by a finite sum of \(L^2\)-norms of right
derivatives of order at most \(d\). Derivatives commute with
the realization of the row, whose Hilbert norm is fixed by its
Petersson normalization. Boundedness of the column on this compact
set and the finite sum thus give
\[
 |B_{t,j}(\varphi,G;y,v)|
    \le C_{t,j,\varphi,G,y}
       \sum_{|\alpha|\le d}\|dM_j(X^{\alpha})v\|_2 .
\]
Here \(X^\alpha\) ranges over ordered monomials in a real Lie
algebra basis. Fixing all but one tensor factor gives the asserted
separate continuity. Compact convolution also proves the density
used when the two-step identity is extended to smooth rows.
\end{proof}

\begin{lemma}\label{lem:joint-period}
Use the word and tuples of Lemma~\ref{hyp:joint-period-towers}. Use the scalar
identity of Proposition~\ref{prop:joint-scalar-descent} with the
same real actions and frames throughout the tuples. Lemma~\ref{lem:word-real-tests-period} identifies the real images and the
continuous periods at each stage. Given finitely
many tuples \(t_i\), coefficients \(\alpha_i\), columns
\(\varphi_i\) of the selected type and terminal rows \(G_i\), one
has the implication
\begin{align*}
 \sum_i\alpha_i h_{t_i} b_{t_i}(U_{t_i}\varphi_i,G_i)=0
 \quad\Longrightarrow\quad
  \sum_i\alpha_i d_{t_i}\,
  B_{t_i,0}\bigl(\varphi_i,D_{t_i,0}\cdots
       D_{t_i,m-1}G_i;y,v\bigr)=0
\end{align*}
for every base frame \(y\) and smooth tensor \(v\).
\end{lemma}

\begin{proof}
Fix \(i\) and write \(t=t_i,\ \alpha=\alpha_i\). At even level \(j\)
use the column \(A_{t,j-2}\cdots A_{t,0}\varphi_i\) and the row
\(D_{t,j}\cdots D_{t,m-1}G_i\), with empty products equal
to the identity. Give this term coefficient
\[
 a_{i,j}=\alpha
 \prod_{0\le b<j/2}(v_{t,2b}k_{t,2b}k_{t,2b+1})
 \prod_{n=j}^{m-1}i_{t,n}.
\]
Thus \(p_ta_{i,m}=\alpha h_t,\ a_{i,0}=\alpha d_t\). Consecutive
coefficients satisfy
\[
 a_{i,j}v_{t,j}k_{t,j}k_{t,j+1}
    =a_{i,j+2}i_{t,j}i_{t,j+1}.
\]

Start with the given relation at \(m\), where the compact pairing
is \(p_tb_t\) times the common real branching map. Assume the
relation at \(j+2\) on every frame and every smooth tensor.
Test at \(j\) on \(I_j(P_1,I_{j+1}(P_2,v'))\).
Equation~\eqref{eq:down-theta} twice expresses the lower period
times \(i_{t,j}i_{t,j+1}\) as a double theta integral. Unfolding
and \eqref{eq:uniform-frame} give a finite sum of upper periods
times \(v_{t,j}k_{t,j}k_{t,j+1}\), integrated on the remaining
archimedean quotients. The frames, the real translates of \(v'\), and their expansion coefficients are common to all indices.

For each such frame and tensor, let \(P_{i,j+2}\) be the corresponding upper period for index \(i\). The induction hypothesis applies pointwise:
\[
 \sum_i a_{i,j+2}P_{i,j+2}=0 .
\]
By the coefficient identity this is the upper-period relation that appears in the unfolded integral. Integration and summation over the common frames both preserve its vanishing.

Two applications of Lemma~\ref{lem:word-real-tests-period} give the relation on \(K_\infty\)-finite
pure tensors. Fix the frame and all factors but one, and denote the resulting
weighted period by \(\mathcal P\). For smooth-topology approximants
\(v_n\) that are \(K_\infty\)-finite, the Sobolev bound on compact
translates gives a continuous seminorm \(p\) with
\[
 |\mathcal P(v)-\mathcal P(v_n)|\le C p(v-v_n)\longrightarrow0 .
\] In each factor
\(K_\infty\)-finite vectors are dense. Extend by continuity one
factor at a time, and descend through \(j=m-2,m-4,\ldots,0\).
\end{proof}

\paragraph{Comparison of two tuples.}

Apply the descended relation to two tuples. To convert that relation into a comparison of their periods, retain the norm ratios of the terminal row and column.

Fix the column similarity \(S\) and its nonzero comparison scalar
\(s_{\rm col}\) for this word. Both are common to the tuples. Apply \(S^{-1}\)
to the initial columns. Lift a specified reference row through the
surjective descent at tuple \(0\) and keep this lift fixed. For the
second input lift its matched row at tuple \(t\), which may itself
be \(0\). The scalars from \eqref{eq:top-bottom} give
\[
   \eta_t=\frac{h_t}{d_t},\qquad c_t=s_{\rm col}\eta_t
   \quad\text{for every tuple }t,\qquad
   \frac{c_t}{c_0}=\frac{h_t d_0}{d_t h_0}.
\]
The terminal reference pairing is evaluated directly on the fixed
lift. The second terminal pairing is read in the comparison
coordinates, where the column map contributes the factor \(s_{\rm col}\).
This distinction persists when the second tuple is \(0\).

\begin{lemma}[Transfer to the bottom pairing]\label{lem:word-two-tuple-transfer}
Let \(A_{\rm ref}\) be the raw finite pairing on the fixed reference
lift and \(B_{\rm ref}\) its bottom pairing. For an input at tuple
\(t\), let \(A_t\) be its finite entry in comparison coordinates
and \(B_t\) its bottom pairing. Then, for \(\alpha,\beta\in\mathbb C\),
\[
 \alpha\eta_0 A_{\rm ref}+\beta c_t A_t=0
       \quad\Longrightarrow\quad \alpha B_{\rm ref}+\beta B_t=0 .
\]
\end{lemma}
\begin{proof}
In Lemma~\ref{lem:joint-period} use the two coefficients
\(\alpha/d_0,\beta/d_t\). On the reference input the top value
before weighting is \(A_{\rm ref}\); on the comparison input the
column similarity and terminal identification give \(s_{\rm col}A_t\).
The top and bottom sums are consequently
\[
 \begin{aligned}
  (\alpha/d_0)h_0 A_{\rm ref}+(\beta/d_t)h_t s_{\rm col} A_t
        &=\alpha\eta_0 A_{\rm ref}+\beta c_t A_t,\\
  (\alpha/d_0)d_0 B_{\rm ref}+(\beta/d_t)d_t B_t
        &=\alpha B_{\rm ref}+\beta B_t .
 \end{aligned}
\]
The descended identity holds for every \(y,v\) and their common
translates. Multiplying by the common conjugate test vector and
integrating in the Poincar\'e formula gives the second equality
as a relation of bottom automorphic pairings. The two translates
have the same level character and use the same column isometry.
\end{proof}

\begin{lemma}[Support of the period ratio]
Let \(S_t\) be the rational primes beneath the twist places.
For some \(N>0\),
\[
 (c_t/c_0)^N\in\mathbb Z[S_t^{-1}]^\times .
\]
For every assigned \(w\mid\ell\), one has \(\ell\notin S_{\mu(w)}\).
\end{lemma}
\begin{proof}
In the quotient of the two products in \eqref{eq:top-bottom},
the fixed level factors and the common masses cancel. In particular,
the common \(s_{\rm col}\) cancels in \(c_t/c_0\), even if \(s_{\rm col}\ne1\).
The stage comparisons are for the ratios of
\(s_{t,j}\), \(k_{t,j}\), \(c_{T,t,j}\), and
\(c_{Q,t,j}/(i_{t,j}i_{t,j+1})\) between the two tuples,
together with the top \(p_t\). Roots of unity from Weil
splittings disappear upon taking powers. The local index and
uniformizer factors left in these ratios are supported on
\(S_t\). Write the remaining factors of \(c_t/c_0\), including any
inverses, as \(u_1,\ldots,u_a\).

For local exponents \(n_j>0\)
with \(u_j^{n_j}\in\mathbb Z[S_t^{-1}]^\times\), take
\(N=\operatorname{lcm}_j n_j\). Equation~\eqref{eq:top-bottom} gives
\[
 (c_t/c_0)^N=\prod_j(u_j^{n_j})^{N/n_j}
                       \in\mathbb Z[S_t^{-1}]^\times.
\] The
row scalars may first be read at a
provisional column system. They are \(s_j\zeta_{t,j}\) with
\(s_j\ne0\) common and \(\zeta_{t,j}\) roots of unity. The
column construction allows these constants as parameters. Its
fixed row data, finite places and level translates agree after
inserting these values, so the same factors occur in the final
calculation.

Fix \(w\mid\ell\) and put \(t=\mu(w)\). At the base tuple the
character is trivial. At a twisted tuple take any host place
\(P\mid\ell\). Its character on inertia is
\[
 \chi_t|_{I_P}=\delta_\ell|_{I_P}\,\psi_w|_{I_P}=1 .
\]
Indeed \(\psi_w\) is unramified at all \(P\mid\ell\). The
support prime of \(\delta_\ell\) is \(p_\delta\) if
\(\ell\ne p_\delta\), and \(p'_\delta\ne p_\delta\) otherwise.
A prime in the extra level of \(t\) must be ramified for
\(\chi_t\), so none lies above \(\ell\). This proves
\(\ell\notin S_{\mu(w)}\) even when several assignments use
the same tuple. \label{cor:period-transport-1}
\end{proof}

\begin{lemma}[Unit at an assigned place]\label{lem:period-unit-at-assigned-place}
Let \(\sigma:L\hookrightarrow\C\) be the selected coefficient
embedding, and let \(w\) lie over a prime outside \(S_t\). If an
identity of period entries has the form
\(\sigma(b)=(c_t/c_0)\sigma(b')\), with \(b,b'\in L\), then their valuations agree whenever the entries are nonzero. More precisely,
\[
 b=b'=0\quad\hbox{or}\quad bb'\ne0,\qquad b/b'\in{\cal O}_{L,w}^{\times}\ \ (bb'\ne0).
\]

\end{lemma}

\begin{proof}
If \(b'=0\), the identity gives \(b=0\). Otherwise set \(a=b/b'\)
and choose \(N>0\) as in the preceding lemma. Writing
\(r=(c_t/c_0)^N\in\mathbb Z[S_t^{-1}]^\times\),
injectivity of \(\sigma\) gives \(a^N=r\), hence
\(N v_w(a)=v_w(r)=0\).
\end{proof}

\begin{lemma}[A nonzero reference entry]\label{lem:nonzero-word-reference-entry}
Suppose the period matrix \(P\) on a \(d\)-dimensional column space,
\(d>0\), satisfies \(\|P^*P-1\|_{\rm HS}^2\le ed\).
If the ramified columns have dimension at least \((1-\delta)d\)
and \(e<1-\delta\), one can choose a reference column and row
with bottom pairing \(b_*\ne0\). Let \(A_*\) be the raw finite
pairing on their fixed lift and put \(p_*=(h_0/d_0)A_*\).
For every finite comparison entry \(d'\) at tuple \(t\) and its
matched bottom pairing \(b'\),
\[
    p_*\ne0,\qquad p_*b'=c_t d' b_* .
\]
\end{lemma}
\begin{proof}
On \(\ker P\) the operator \(P^*P-1\) is minus the identity. Extending
an orthonormal basis of the kernel to the whole column space gives
\[
 \dim\ker P\le\|P^*P-1\|_{\rm HS}^2\le ed<(1-\delta)d.
\]
Choose a ramified column outside that kernel, a row on which it
pairs to \(b_*\ne0\), and a preimage under the column similarity.
Apply Lemma~\ref{lem:word-two-tuple-transfer} with
\((\alpha,\beta)=(1,0)\). If \(p_*=0\), its conclusion forces
\(b_*=0\). Thus both \(p_*\) and \(A_*\) are nonzero.

For the fixed reference entry and any comparison entry \(d'\) take
\((\alpha,\beta)=(c_t d',-p_*)\) in the same lemma. The finite
relation and the resulting bottom relation are
\[
 \begin{aligned}
  (c_t d')\eta_0 A_*-p_*c_t d'&=0,\\
  (c_t d')b_*-p_*b'&=0 .
 \end{aligned}
\]
Restoring the raw finite reference entry gives
\begin{equation}\label{eq:raw-word-reference}
       (h_0/d_0)A_* b'=c_t d' b_* .
\end{equation}
In particular \(d'=0\)
forces \(b'=0\).
\end{proof}

To isolate the scalar used at finite places, put \(K_u=b_*/p_*\)
for the chosen reference entry of the word \(u\). Apply
Lemma~\ref{lem:nonzero-word-reference-entry} twice, keeping that entry
fixed and using comparison inputs at \(t\) and at \(0\), respectively.
For their finite and bottom entries \((A_t,B_t)\), \((A_0,B_0)\) this gives
\begin{equation}\label{eq:word-period-cross-product}
 B_t=K_u c_t A_t,\qquad B_0=K_u c_0 A_0,\qquad
 B_t A_0=(c_t/c_0)B_0 A_t .
\end{equation}
The last equality also holds when either comparison entry vanishes.
Dividing the column similarity by its positive norm ratio gives a column
isometry. Central right translation gives a common row isometry. These
operations use the same scalars at \(t\) and \(0\). Thus the tuple
ratio in Lemma~\ref{lem:period-unit-at-assigned-place} has the indicated
orientation after passing to the isometric coordinates.

\subsection{The integral condition on this comparison}

Descent identifies the complex periods. A divisibility statement requires identifying the integral rows inside that common complex space and retaining the metric at every embedding. First identify the
rational lattice whose image in the complex row space occurs in the period, then
check the coefficient lattice on the negative line.

\paragraph{Rational and integral row coordinates.}
In the rational realization of a row let \(L_0\)
be the \(\mathcal O_L\)-lattice in \(L^d\), \(H\) its Gram
matrix at \(\sigma\), and \(E_0\) the rational matrix of the
good-part projector. Use the same construction at each other
infinite place with its specified Gram matrix. In the selected
complex row space choose the embedding \(j\) on columns
satisfying
\[
 \langle jx,jy\rangle=\sigma(x^*Hy),\qquad
 P_{\rm good}\,jx=j(E_0x),\qquad
           \operatorname{span}_{\C} jL^d=W .
\]


These identifications preserve every local determinant index once the integral generators on both sides have been matched. Indeed, if \(M,M'\) are two local lattices with basis matrices \(A,B\) and the common coordinate change is \(U\), then
\[
 v_w\det((UA)^{-1}(UB))=v_w\det(A^{-1}B).
\]
For a coefficient-field linear isometry \(J\) of the rational Hermitian spaces, \(\langle Jx,JM\rangle=\langle x,M\rangle\) over the coefficient field. Requiring these pairings to lie in its integer ring gives \((JM)^\#=J(M^\#)\). This last calculation uses the rational identification with the integral generators at all places. The first displayed metric identity forces injectivity.
In coordinates it is realized by a square root of the positive
Hermitian matrix \(\sigma H\), followed by an isometry
to \(W\). The identification of the projector is a further
condition verified on a basis in the automorphic model.
\begin{lemma}[Integral row and negative-line lattices]
If \(g_1,\ldots,g_l\) generate \(L_0\), then

\[
 j^{-1}\!\left(\sum_i\sigma(\mathcal O_L)jg_i\right)=L_0.
\]
 If the integral
Witt identification sends the descent line to \(\mu z_u\), where
\(z_u\) is the primitive last vector of squared norm \(-T_u\),
its coefficient lattice and squared norm are respectively
\(\mu^{-1}\mathcal O_{E_{\rm CM}}\) and \(-\mu\bar\mu T_u\).
\end{lemma}
\begin{proof}
The integral row in these coordinates is generated over \(\sigma(\mathcal O_L)\)
by \(jg_i\). Consequently its inverse image under \(j\) is
exactly \(L_0\). Subtract the same integral combination in
\(L^d\) and use injectivity. This works for projective
lattices with any finite generating set. The rational
identification with the division lattice requires equality of
this module and of the Hermitian metrics at \emph{all} infinite
places. In particular, an arbitrary rescaling of a complex
model can change the application to slopes.

For the integral primitive flag at word \(u\), the class of \(z_u\)
in the orthogonal quotient generates a primitive rank-one summand.
Thus \(a\mu z_u\) integral implies \(a\mu\in\mathcal O_{E_{\rm CM}}\)
by passage to this summand. The converse follows from integrality
of \(z_u\). Hermitian homogeneity gives the asserted squared norm.
\end{proof}

\begin{lemma}[Comparison on the rational good part]
Write \(B_{\rm ch}\) and \(B_{\rm per}\) for the pairings on the selected base rows in chart and period coordinates, respectively. On the image of \(E_0\), evaluation on pairs of basis vectors determines the matrix after scalar extension. In corresponding field bases \(e_i,f_k\) it is the entrywise identity
\[
 B_{\rm ch}(e_i,f_k)=B_{\rm per}(e_i,f_k)\qquad\text{for all }i,k.
\]

\end{lemma}
\begin{proof}
The row test functions and arithmetic maps are compared
on the exact range of \(E_0\). The tensor evaluation used for
the automorphic realization is injective on the base row model.
On the field basis the two realizations give the same
evaluations by the chart identities. They therefore give the
same images on this range after scalar extension. Schur
uniqueness on the irreducible archimedean type allows one scalar
in the conjugate evaluation; it cancels in the similarity ratio.
For field bases \(e_i,f_k\) of the two ranges, every complex
bilinear extension satisfies
\[
 B\left(\sum_i x_ie_i,\sum_k y_k f_k\right)=\sum_{i,k}x_iy_k B(e_i,f_k).
\]
Hence the basis identities give equality on all pairs. The map
onto the selected base rows is surjective, yielding their period
matrices. The fixed-level support used in
the comparison is proved in Section~\ref{subsec:fixed-level-support}.
The arithmetic application uses the integral frame and tuple
identifications of Section~\ref{subsec:division-level-pins},
and their geometric and division bounds at the selected weights.
\end{proof}
